\documentclass[conference]{IEEEtran}
\IEEEoverridecommandlockouts
% The preceding line is only needed to identify funding in the first footnote. If that is unneeded, please comment it out.
\usepackage[T1]{fontenc}
\usepackage{cite}
\usepackage{amsmath,amssymb,amsfonts}
\usepackage{graphicx}
\usepackage{booktabs}
\usepackage{url}
\usepackage{cleveref}   % keep this LAST

\def\BibTeX{{\rm B\kern-.05em{\sc i\kern-.025em b}\kern-.08em
    T\kern-.1667em\lower.7ex\hbox{E}\kern-.125emX}}
\begin{document}

\title{Model Capacity or Input Resolution? Density-Based Counting of
\textit{Tenebrio molitor} Larvae under Edge Memory Constraints\\
\thanks{This thesis was prepared in partial fulfilment of the requirements for the Degree of Bachelor of Computer Science, Maastricht University. Supervisor: Petra Bosilj}
}

\author{\IEEEauthorblockN{Adriane Bien}
\IEEEauthorblockA{\textit{Department of Advanced Computing Sciences} \\
\textit{Faculty of Science and Engineering}\\
\textit{Maastricht University}\\
Maastricht, The Netherlands}
}

\maketitle

\begin{abstract}
Commercial insect farming depends on automated monitoring of rearing trays,
but existing machine-vision work for this task relies on detection and
segmentation, which degrade when larvae are densely packed and occluded.
Density-based crowd counting suits such scenes yet remains unexplored in insect
agriculture. This thesis applies it to Tenebrio molitor larvae under edge memory
constraints, asking how a fixed budget should be divided between model capacity
and input resolution. A heavyweight network, CSRNet, and a lightweight one,
MobileCount, are tuned independently on the TenebrioVision dataset and evaluated
across a six-step resolution ladder, with peak inference memory recorded for
each of the twelve resulting configurations. Both count to within one larva at a
sixteenth of the native resolution. Peak memory is weight-dominated at the three
coarsest steps and activation-dominated above them, so downscaling stops paying
well before accuracy stops improving, and MobileCount holds every point on the
accuracy-memory frontier except its high-accuracy corner. On a Jetson Nano all
twelve configurations fit, so the required accuracy rather than the footprint
sets the operating point.
\end{abstract}

\begin{IEEEkeywords}
crowd counting, density estimation, edge computing, insect farming,
precision livestock farming, convolutional neural networks
\end{IEEEkeywords}

\section{Introduction}
To meet the growing global demand for protein, insect farming has emerged as a
promising alternative to conventional livestock production~\cite{vanhuis2013}.
\textit{Tenebrio molitor}, the yellow mealworm, is one of the few insect species
approved for human consumption in Europe~\cite{eu2021tenebrio}. To be
commercially viable, however, production must scale up significantly---a challenge
that hinges directly on automation~\cite{rumpold2013}.

Automated mealworm farming requires continuous monitoring of the rearing plates
to inform decisions about feeding, harvesting, and environmental control. Prior work has explored machine vision for this purpose, but most of it relies on object detection and instance segmentation~\cite{majewski2022multipurpose}. Papadopoulos et al.~\cite{Papadopoulos2024TenebrioVision} found that these
paradigms degrade on real-farm images of \textit{T.\ molitor} rearing trays, where larvae are densely packed and heavily occluded. Crowd counting via density estimation offers a far more suitable framework for these
conditions~\cite{Sindagi2017ASO}, yet it remains entirely unexplored in insect
agriculture.

The primary driver for automation is economic: reducing operational costs and
maximising production efficiency~\cite{rumpold2013}. Because the industry lacks
standardization~\cite{NAWOYA2024108503}, farms have adopted a diverse range of
technological solutions~\cite{CoRoSect2022AdvancedArchitecture,
kroncke2020automation}. One compelling approach is edge AI, which is
particularly attractive for cost-sensitive operations because it runs inference
locally on low-cost embedded devices~\cite{angulo2025realtimecricketsortingsex}.
This thesis therefore investigates how crowd-counting models can be deployed
within the strict memory constraints of edge hardware, specifically the Jetson Nano. 

Memory consumption during inference is governed primarily by two factors:
larger architectures need more memory for their parameters, and
higher-resolution inputs enlarge the intermediate feature maps held during the
forward pass~\cite{Chao2019HarDNet}. Reducing the footprint therefore comes down
to shrinking the model or shrinking the input.

This thesis first establishes that density-based counting is viable in this
setting, then asks:

\begin{quote}
Under the memory constraints of edge hardware, how should a fixed memory budget
for density-based mealworm counting be allocated between model capacity and
input resolution, and which resulting configuration best serves automated
rearing-tray monitoring?
\end{quote}

\noindent This question is decomposed into four sub-questions:

\begin{itemize}\setlength{\itemsep}{0pt}\setlength{\parskip}{0pt}
  \item \textbf{RQ1:} How does input resolution affect the counting accuracy of
    differently sized crowd-counting models?
  \item \textbf{RQ2:} How does peak inference memory scale with input resolution
    and with the size of the backbone?
  \item \textbf{RQ3:} Which architecture--resolution combinations are
    Pareto-optimal on the accuracy--memory plane?
  \item \textbf{RQ4:} Which configurations fit the Jetson Nano's budget while
    meeting the accuracy demanded by operational insect rearing?
\end{itemize}

\section{Related Work}
\subsection{Crowd Counting Approaches}

Crowd counting methods can be broadly categorised into three main approaches~\cite{gao2020crowdcountingsurvey, li2018csrnet}: detection-based~\cite{dalal2005histograms, felzenszwalb2010object}, regression-based~\cite{chan2008privacy, idrees2013multi}, and density estimation-based~\cite{lempitsky2010learning} methods.

The earliest approaches in the field were detection-based, which produce a population estimate by aggregating individual bounding-box predictions. Typical examples include pedestrian detectors based on hand-crafted Histogram of Oriented Gradients (HOG) features~\cite{dalal2005histograms} and part-based models that detect individual body parts to better handle partial occlusion~\cite{felzenszwalb2010object}. These methods struggle under severe occlusion and high object density~\cite{gao2020crowdcountingsurvey}.

Regression-based methods were developed to address this by computing object counts directly from image features, bypassing the need for explicit localization~\cite{chan2008privacy, idrees2013multi, gao2020crowdcountingsurvey}. While this shift makes them significantly more robust to heavy occlusion, they sacrifice spatial information entirely, leaving no insight into how individuals are actually distributed across a scene~\cite{gao2020crowdcountingsurvey}.

The most widely used crowd-counting models today rely on density estimation~\cite{lempitsky2010learning, gao2020crowdcountingsurvey}. Rather than localising individuals explicitly, this approach learns to map an input image to a continuous density map whose integral over any region gives the count within it, and over the whole map the total. Because the map records where mass is concentrated, it retains the spatial information that regression discards, and such models are typically trained by regressing the predicted map against a generated ground-truth one~\cite{wang2020mobilecount, li2018csrnet, gao2020crowdcountingsurvey}.

\subsection{CNN-Based Density Estimation}

Convolutional Neural Networks (CNNs) learn hierarchical feature representations directly from raw pixel data rather than relying on hand-crafted features and shallow regressors~\cite{krizhevsky2012imagenet}, which buys far greater robustness to occlusion, lighting variation and background clutter and has made them the de facto standard for density estimation~\cite{gao2020crowdcountingsurvey, li2018csrnet}.

Early approaches used shallow, single-branch networks of standard convolutional and fully connected layers that regress a count or coarse density map directly~\cite{fu2015fast, wang2015deep}, but these struggled with the extreme scale variation inherent to crowded scenes~\cite{Zhang_2016_CVPR, gao2020crowdcountingsurvey}. That limitation motivated the development of multi-column architectures such as the Multi-column Convolutional Neural Network (MCNN), which runs parallel network branches with different receptive-field sizes to capture image features at multiple scales simultaneously~\cite{Zhang_2016_CVPR}.

The field has since largely shifted back toward single-column designs~\cite{gao2020crowdcountingsurvey}---networks with a single feature-extraction path, as in the early basic CNNs, but now substantially deeper and built on pretrained classification backbones rather than shallow stacks of plain layers~\cite{li2018csrnet, gao2020crowdcountingsurvey}. This design has become one of the most widely adopted frameworks in crowd counting~\cite{gao2020crowdcountingsurvey}; the Congested Scene Recognition Network (CSRNet)~\cite{li2018csrnet} served as a pivotal model in this transition and consequently remains a standard baseline for benchmarking newer methods, particularly in highly dense scenes.

\subsection{Lightweight Networks}
Beyond core architectural designs, lightweight models have recently emerged as a crucial area of research, driven by the need for resource-constrained deployment on edge devices~\cite{wang2020mobilecount}. The primary objective here is to minimise computational overhead and parameter counts without incurring a significant penalty to overall counting accuracy~\cite{gao2020crowdcountingsurvey, wang2020mobilecount}.

A prominent example of this in density-based crowd counting is MobileCount~\cite{wang2020mobilecount}, an encoder-decoder network built on top of a lightweight MobileNetV2 backbone~\cite{sandler2018mobilenetv2}. Compared directly with CSRNet, MobileCount reduces the overall model footprint by roughly 80\%---requiring only 3.40 million parameters versus CSRNet's 16.26 million---while requiring 52 times fewer floating-point operations (FLOPs) and exhibiting only a minor reduction in accuracy across standard benchmark datasets~\cite{wang2020mobilecount}.

\section{Methodology}
\subsection{Models}

Two CNN-based density-estimation models are trained and evaluated: the accurate but computationally heavy CSRNet~\cite{li2018csrnet}, and the lightweight, efficiency-oriented MobileCount~\cite{wang2020mobilecount}. Together, they represent complementary points on the accuracy--efficiency trade-off.

\subsubsection{CSRNet}
CSRNet pairs a front-end feature extractor with a dilated convolutional back-end~\cite{yu2016multiscale, li2018csrnet}. The front-end comprises the first ten layers of a pre-trained VGG-16 network~\cite{simonyan2015vgg} with the classification layers removed; its three max-pooling operations reduce the spatial resolution to $H/8 \times W/8$ while providing strong transfer-learning capacity. Rather than downsampling further, the back-end uses dilated convolutions to enlarge the receptive field without reducing resolution or adding parameters, aggregating multi-scale context at a constant spatial size~\cite{yu2016multiscale, li2018csrnet}. Configuration B is used here, applying a uniform dilation rate of 2 throughout the back-end and identified in the original paper as the most accurate of the four proposed~\cite{li2018csrnet}. A final $1 \times 1$ convolution collapses the feature maps to a single channel at $H/8 \times W/8$, bilinearly upsampled by a factor of eight to the input resolution.

\subsubsection{MobileCount}
MobileCount is an encoder--decoder network~\cite{wang2020mobilecount}. Its encoder is based on MobileNetV2~\cite{sandler2018mobilenetv2}, whose inverted residual blocks with linear bottlenecks separate spatial filtering from channel mixing through depthwise convolutions, substantially reducing the parameter count relative to standard convolutions. MobileCount retains only four of MobileNetV2's seven bottleneck blocks and inserts an additional stride-2 max-pooling layer at the start of the encoder, both of which the original paper reports cut FLOPs considerably at a tolerable cost to accuracy~\cite{wang2020mobilecount}. The decoder follows the Light-Weight RefineNet~\cite{nekrasov2018lightweight}, progressively fusing the four encoder feature maps---spanning $H/4 \times W/4$ to $H/32 \times W/32$---through chained residual pooling and fusion blocks back up to $H/4$. A final $1 \times 1$ convolution maps this to a single channel, bilinearly upsampled by a factor of four to the input resolution~\cite{wang2020mobilecount}.

\subsection{Ground Truth Generation}

Point annotations are derived from the centre of each bounding box to construct a discrete point map, which is then convolved with a Gaussian kernel to give the ground-truth density map~\cite{lempitsky2010learning, Zhang_2016_CVPR}. Let $\Omega$ denote the pixel grid of a training image, so that a density map assigns a scalar to every pixel $\mathbf{x} \in \Omega$. For image $i$ containing $M_i$ annotated larvae at positions $\mathbf{x}_1, \dots, \mathbf{x}_{M_i}$, the ground truth $D_i^{GT}$ is defined pixelwise as:
\begin{equation}
  \label{eq:density-map}
  D_i^{GT}(\mathbf{x}) = \sum_{j=1}^{M_i} \left( \delta(\mathbf{x} - \mathbf{x}_j) * G_{\sigma} \right)(\mathbf{x}),
  \qquad \mathbf{x} \in \Omega
\end{equation}

where $\delta(\mathbf{x} - \mathbf{x}_j)$ is a Dirac delta function centred at the larva position $\mathbf{x}_j$ and $G_{\sigma}$ is an isotropic Gaussian kernel with standard deviation $\sigma$. The kernel is normalized to unit mass, so that summing $D_i^{GT}$ over all pixels recovers the annotated count $M_i$.

\subsection{Loss Function}

Both models produce a predicted density map $D_i^{PM}$ on the same grid $\Omega$ as the ground truth. Training minimises the mean squared error between the two, computed pixel by pixel and averaged over all pixels and all images in a mini-batch~\cite{li2018csrnet, wang2020mobilecount}:
\begin{equation}
  \label{eq:mse-loss}
  \mathcal{L} = \frac{1}{N \lvert \Omega \rvert} \sum_{i=1}^{N} \sum_{\mathbf{x} \in \Omega}
  \left( D_i^{PM}(\mathbf{x}) - D_i^{GT}(\mathbf{x}) \right)^2
\end{equation}
where $N$ is the number of images in the mini-batch and $\lvert \Omega \rvert$ is the number of pixels per image. The population estimate for an image is then obtained by summing the predicted density map over the full pixel grid~\cite{lempitsky2010learning}.

\section{Experiments}
\subsection{Dataset}
This study uses the TenebrioVision dataset~\cite{Papadopoulos2024TenebrioVision},
comprising 1{,}120 images at $3088 \times 2076$ pixels, grouped into 14 count
levels of 80 images each. Every image in a level contains a fixed number of
\textit{Tenebrio molitor} larvae---10, 15, 20, 25, 30, 35, 40, 45, 50, 60, 70,
80, 90 or 100---and each larva is labelled with a bounding box, giving 53{,}600
annotated instances. It is the only publicly available dataset reflecting the
high-density conditions of industrial larvae farming.

Following the split protocol of~\cite{Papadopoulos2024TenebrioVision}, the data
is stratified by count level and divided into training, validation and test sets
in a $75:15:10$ ratio, shuffling within each level with a fixed seed (42). This
yields 840 training, 168 validation and 112 test images.

All experiments operate on a resolution ladder derived from this data by
progressively halving the original images (rounding up on odd dimensions),
giving $1544\times1038$, $772\times519$, $386\times260$, $193\times130$,
$97\times65$ and $49\times33$. Resizing is applied identically to all three
subsets, preserving the split across variants.

\subsection{Implementation Details}
\label{sec:impl_details}

All experiments are implemented in the C\textsuperscript{3}
Framework~\cite{gao2019c3}, whose authors tuned CSRNet and MobileCount on
ShanghaiTech Part~B. That dataset shares the fixed viewpoint, roughly constant
object size and comparatively uniform density of TenebrioVision, so its
published configuration is used as the baseline
(Table~\ref{tab:baseline_hyperparameters}). Weight decay is fixed at the
framework default of $1\times10^{-4}$ and the learning rate follows the default
step decay of $\gamma = 0.995$ per epoch. All training, evaluation and
measurement is performed in PyTorch on an NVIDIA GeForce RTX 5070~Ti with
16\,GB of VRAM.

\begin{table}[t]
    \centering
    \caption{C\textsuperscript{3} Framework (baseline) hyperparameters.}
    \label{tab:baseline_hyperparameters}
    \begin{tabular}{lcc}
        \toprule
        Hyperparameter          & CSRNet             & MobileCount \\
        \midrule
        Optimiser               & \multicolumn{2}{c}{Adam} \\
        Weight decay            & \multicolumn{2}{c}{$1\times10^{-4}$} \\
        Batch size              & \multicolumn{2}{c}{$6$} \\
        LR decay                & \multicolumn{2}{c}{step, $\gamma = 0.995$} \\
        Initial learning rate   & $1\times10^{-5}$   & $1\times10^{-4}$ \\
        Density-target scale $k$ & $100$             & $2550$ \\
        \bottomrule
    \end{tabular}
\end{table}

\paragraph{Density-target scaling.}
The framework multiplies the density-map regression target by a constant $k$ and
divides the prediction by $k$ at evaluation. Since a density map sums to the
object count, its per-pixel magnitude is inversely proportional to pixel count,
making $k$ resolution-dependent. The published values ($k = 100$ for CSRNet,
$k = 2550$ for MobileCount) are anchored to $772\times519$, the variant closest
in pixel count to ShanghaiTech Part~B, and scaled in proportion elsewhere:
\begin{equation}
    k(r) \;=\; k_{\mathrm{base}}\cdot \frac{A(r)}{A(772\times519)},
    \label{eq:logpara_scaling}
\end{equation}
where $A(\cdot)$ denotes pixel count. Here $k$ is a deterministic function of
resolution rather than a tuned quantity.

\paragraph{Ground-truth generation.}
Ground truth is generated via Eq.~\eqref{eq:density-map} using the fixed-kernel
strategy CSRNet
applies on ShanghaiTech Part~B, giving an object-to-kernel ratio of
$\sigma = \mathrm{bbox}/2$, where $\mathrm{bbox}$ is the average larva
bounding-box size in TenebrioVision. Scaled proportionally with resolution, this
yields $\sigma = 24, 12, 6, 3, 2, 1$ pixels from the largest variant to the
smallest, rounded to the nearest pixel at the two smallest scales. The kernel is
truncated to an effective window of $8\sigma + 1$ pixels.

\subsection{Experimental Design}

The evaluation proceeds in two stages. Stage~1 establishes a training
configuration for each architecture; Stage~2 holds that configuration fixed and
varies only the input resolution. No hyperparameter is tuned per resolution.

\subsubsection{Stage 1: Configuration experiments}

Tuning is performed independently for each architecture. Both models start from
the baseline configuration of Table~\ref{tab:baseline_hyperparameters}, and each
experiment sweeps one parameter while the others are held at their baseline
values. The value determined in each experiment is used in all subsequent
experiments.

\paragraph{Tuning resolution and budget.}
Every tuning run is trained for a fixed budget of 1500 epochs, by which point
the per-epoch decay of $\gamma = 0.995$ has reduced the learning rate to roughly
$5\times10^{-4}$ of its initial value. Tuning is performed at $386\times260$,
the mid-range point of the ladder, where a run takes ${\approx}1$\,h for
MobileCount and ${\approx}2.7$\,h for CSRNet; at $1544\times1038$ a single
CSRNet run takes ${\approx}26$\,h.

\paragraph{E1 --- Augmentation scheme.}
Three schemes are compared on both architectures: no augmentation; the online
scheme of MobileCount, a random crop spanning $80\%$ of the image drawn afresh
at every iteration~\cite{wang2020mobilecount}; and the offline scheme of CSRNet,
which expands each image into $18$ fixed patches, namely four non-overlapping
quarters and five random quarter-size crops, each duplicated by horizontal
mirroring~\cite{li2018csrnet}.

\paragraph{E2 --- Optimiser and learning rate.}
Since weight decay is non-zero, Adam is compared with AdamW (decoupled weight
decay) across a shared learning-rate grid, spanning $10^{-6}$--$10^{-3}$ for
CSRNet and $10^{-5}$--$10^{-2}$ for MobileCount.

\paragraph{E3 --- Batch size.}
The candidates $B\in\{1,2,4,6,8\}$ include $B=1$, used in the original CSRNet
work~\cite{li2018csrnet}, and $B=6$, the framework default. Because the optimal
learning rate for adaptive optimisers scales with the square root of the batch
size~\cite{malladi2022sdes}, each candidate is evaluated at
\begin{equation}
    \eta(B) \;=\; \eta^\star\sqrt{B/6},
    \label{eq:lr_batch_scaling}
\end{equation}
anchored at $B=6$, the batch size held fixed while $\eta^\star$ was selected in
E2, so the anchor cell reproduces the E2 configuration exactly.

\paragraph{E4 --- Epoch budget across resolutions.}
The configuration selected in E1--E3 is run once at every resolution for the
full 1500-epoch budget to locate where convergence occurs along the ladder.
This experiment selects no hyperparameter; it fixes the training budget carried
into Stage~2.

\subsubsection{Stage 2: Cross-resolution comparison}
Each architecture's Stage~1 configuration is held fixed while the input
resolution is varied across the six ladder variants. Each of the twelve
architecture--resolution cells is trained, selected and evaluated over five
independent runs, and results are reported as the aggregate across those runs.

\subsection{Evaluation Metrics}
\label{sec:eval_metrics}
Following standard practice in the crowd-counting
literature~\cite{gao2020crowdcountingsurvey, li2018csrnet}, accuracy is measured
by Mean Absolute Error (MAE) and Mean Squared Error (MSE):

\begin{equation}
    \text{MAE} = \frac{1}{N} \sum_{i=1}^{N} \left| y_i - \hat{y}_i \right|,
    \label{eq:mae}
\end{equation}
\begin{equation}
    \text{MSE} = \sqrt{\frac{1}{N} \sum_{i=1}^{N}
    \left( y_i - \hat{y}_i \right)^2},
    \label{eq:mse}
\end{equation}
where $N$ is the number of test images, $y_i$ the ground-truth count and
$\hat{y}_i$ the predicted count. Following the convention of the crowd-counting
literature, MSE here denotes the \emph{root} mean squared error, so that both
metrics are expressed in larvae and are directly comparable; it penalises large
errors more heavily and is therefore sensitive to outliers.

Efficiency is measured by peak inference memory, defined as the PyTorch
allocator's peak memory consumption during a single forward pass, covering model
weights, the input tensor and the maximum intermediate activations held
simultaneously. Because this is a deterministic property of a given
configuration, it is measured once per experimental cell.

\section{Results}
\label{sec:results}

Throughout, italics mark the C\textsuperscript{3} Framework default where one
was among the evaluated settings and bold marks the best result in each block.
Unless stated otherwise, tuning runs use $386\times260$ and selection is on
validation MAE.

\subsection{Hyperparameter Tuning}
\label{subsec:res_tuning}

This subsection reports the three sequential One-Factor-at-a-Time (OFAT) tuning
stages --- augmentation, optimiser and learning rate, and batch size --- and
closes with the configuration carried into the cross-resolution experiments.

\begin{table}[t]
    \centering
    \caption{E1 --- Augmentation tuning --- Training- and validation-split MAE at the
    best-validation checkpoint.}
    \label{tab:aug_paper_schemes}
    \small
    \begin{tabular}{l rrr rrr}
        \toprule
        & \multicolumn{3}{c}{CSRNet} & \multicolumn{3}{c}{MobileCount} \\
        \cmidrule(lr){2-4} \cmidrule(lr){5-7}
        Scheme        & val & test & epoch & val & test & epoch \\
        \midrule
        none          & $3.514$ & $3.280$          & $8$   & $0.605$ & $0.760$          & $150$ \\
        $80\%$ crop   & $3.513$ & $\mathbf{3.097}$ & $14$  & $0.632$ & $\mathbf{0.618}$ & $185$ \\
        $18$-patch    & $4.074$ & $3.549$          & $640$ & $0.622$ & $0.700$          & $112$ \\
        \bottomrule
    \end{tabular}
\end{table}

Table~\ref{tab:aug_paper_schemes} compares the three augmentation schemes. The
online crop attains the lowest test MAE for both architectures and the offline
$18$-patch scheme the highest. The validation split does not reproduce this
ordering: for CSRNet the crop and no-augmentation runs are separated by $0.001$,
a single unit in the last reported digit, and for MobileCount the
no-augmentation run is nominally the best of the three.

\begin{table}[t]
    \centering
    \caption{E2 --- Best validation MAE over the optimiser $\times$ learning-rate
    grid.}
    \label{tab:lr_grid}
    \begin{tabular}{ccc@{\hspace{2em}}ccc}
        \toprule
        \multicolumn{3}{c}{CSRNet} & \multicolumn{3}{c}{MobileCount}\\
          LR& \emph{Adam} & AdamW & LR & \emph{Adam} & AdamW \\
        \midrule
        $10^{-6}$ & $2.74$          & $0.923$          & $10^{-5}$ & $0.880$          & $0.831$ \\
        $10^{-5}$ & $\mathit{3.51}$ & $0.394$          & $10^{-4}$ & $\mathit{0.632}$ & $0.542$ \\
        $10^{-4}$ & $4.33$          & $\mathbf{0.320}$ & $10^{-3}$ & $0.420$          & $\mathbf{0.385}$ \\
    \end{tabular}
\end{table}

Table~\ref{tab:lr_grid} shows AdamW yielding lower validation MAE at every grid
point, with each architecture's optimum at the highest rate at which training
remained stable. The top decade of each sweep --- $10^{-3}$ for CSRNet and
$10^{-2}$ for MobileCount --- is omitted from the table: at those rates neither
network learned, validation MAE plateauing at approximately $22.88$ and $23.14$
respectively --- close to the $23.3$ mean absolute deviation of the count
distribution, and therefore consistent with a collapse to a constant
prediction.

\begin{table}[t]
\centering
\caption{E3 --- Batch-size sweep with the learning rate rescaled
from each architecture's optimal learning rate by $\sqrt{b/6}$.}
\label{tab:batchsize}
\begin{tabular}{l cc@{\hspace{2em}}cc}
\toprule
 & \multicolumn{2}{c}{CSRNet} & \multicolumn{2}{c}{MobileCount} \\
\cmidrule(lr){2-3} \cmidrule(lr){4-5}
Batch size & LR & MAE & LR & MAE \\
\midrule
$1$          & $4.0\times10^{-5}$ & $\mathbf{0.288}$ & $4.0\times10^{-4}$ & $2.85$ \\
$2$          & $6.0\times10^{-5}$ & $0.292$          & $6.0\times10^{-4}$ & $0.543$ \\
$4$          & $8.0\times10^{-5}$ & $0.307$          & $8.0\times10^{-4}$ & $0.436$ \\
$\mathit{6}$ & $1.0\times10^{-4}$ & $0.320$          & $1.0\times10^{-3}$ & $\mathbf{0.385}$ \\
$8$          & $1.2\times10^{-4}$ & $0.319$          & $1.2\times10^{-3}$ & $0.460$ \\
\bottomrule
\end{tabular}
\end{table}

The two architectures respond to batch size in opposite directions
(Table~\ref{tab:batchsize}). CSRNet improves as the batch size falls below the
default and is best at batch size $1$; the trend is monotone over every step
except the last, where the two largest batches differ by $0.001$, a single unit
in the last reported digit. MobileCount instead has an interior optimum at the
default and degrades on both sides of it, collapsing by a factor of seven at
batch size $1$.

E4 runs the selected configuration once at every resolution for the full
$1500$-epoch budget. Across these twelve runs the best-validation epoch falls
between epoch $33$ and epoch $445$, and validation MAE at epoch $1500$ is higher
than at that checkpoint in every case, so the second half of the budget only
overfits. An $800$-epoch budget is therefore carried into Stage~2.
The checkpoint values themselves (Table~\ref{tab:e4_budget}) show CSRNet
improving monotonically with resolution, whereas MobileCount is best at
$386\times260$ and degrades above it.

\begin{table}[t]
    \centering
    \caption{Final selected configurations.}
    \label{tab:final_configs}
    \begin{tabular}{lcc}
        \toprule
        & CSRNet & MobileCount \\
        \midrule
        Optimiser             & \multicolumn{2}{c}{AdamW} \\
        Initial learning rate & $4\times10^{-5}$ & $1\times10^{-3}$ \\
        Batch size            & $1$ & $6$ \\
        Weight decay          & \multicolumn{2}{c}{$1\times10^{-4}$} \\
        LR schedule           & \multicolumn{2}{c}{step $\gamma=0.995$} \\
        Training budget       & \multicolumn{2}{c}{$800$ epochs} \\
        Augmentation          & \multicolumn{2}{c}{online $80\%$ crop} \\
        Density-target scale $k$ & $100$ & $2550$ \\
        & \multicolumn{2}{c}{\footnotesize scaled per Eq.~\eqref{eq:logpara_scaling}} \\
        Random seed      & \multicolumn{2}{c}{$3035$} \\
        \bottomrule
    \end{tabular}
\end{table}

The resulting configuration is summarised in Table~\ref{tab:final_configs}.
These settings are held fixed for the remainder of the evaluation; only the
input resolution and the density-target scale $k$, which is rescaled with
resolution according to Eq.~\eqref{eq:logpara_scaling}, vary across the
cross-resolution experiments.

% ============================================================================
\subsection{Cross-Resolution Evaluation}
\label{subsec:res_resolution}

Table~\ref{tab:resolution_test} reports test error for both architectures
across the resolution ladder, aggregated over five runs;
Figure~\ref{fig:density_maps} shows the corresponding predicted density maps
for a representative test image.
\begin{figure}
    \centering
    \includegraphics[width=1\linewidth]{images/test_error_vs_resolution_repeats.png}
    \caption{Test (a)~MAE and (b)~MSE against input resolution, on log axes.
    Markers are means over five runs and bars span one standard deviation.}
    \label{fig:allres_trend}
\end{figure}
Point estimates of both metrics decrease monotonically with resolution for
both architectures, and both attain their lowest error at $1544\times1038$.
Table~\ref{tab:resolution_test} and Figure~\ref{fig:allres_trend} show how the
two ladders differ in shape. CSRNet reduces MAE by a near-constant factor across
three consecutive doublings ($36\%$, $48\%$, $41\%$), tapering to $15\%$ on the
final step. MobileCount's gains are unevenly distributed: two steps are
effectively flat ($49\times33 \to 97\times65$, $6\%$; $386\times260 \to
772\times519$, $9\%$, both within one standard deviation), separated by two
large ones ($64\%$ and $48\%$) and followed by a further $27\%$ reduction at the
top of the ladder.

The smallest resolution is a special case for CSRNet. Its MAE of $6.55$ is more
than twice MobileCount's, its best checkpoint occurs at epoch $33$ --- the
earliest of any CSRNet cell, and the only one at which all five runs peaked
together --- and its predicted density map carries no spatial structure
(Figure~\ref{fig:density_maps}), the count arising from a near-uniform mass
rather than from localised responses.

The ordering between the architectures reverses once on MAE, and only at the
coarse end. MobileCount is lower at $49\times33$; CSRNet is lower at
$97\times65$; the two are indistinguishable at $193\times130$ and
$386\times260$, where the differences ($0.006$ and $0.001$) are an order of
magnitude smaller than the run-to-run spread; and CSRNet is lower at
$772\times519$ and $1544\times1038$ with non-overlapping intervals.

The two metrics do not agree. MobileCount attains the lower MSE at every
resolution except $97\times65$ --- including $772\times519$ and
$1544\times1038$, where CSRNet's MAE is clearly lower, although at those two the
intervals overlap. The ratio between the metrics separates the architectures
more cleanly than either alone: MobileCount's MSE/MAE ratio is near-constant
across the ladder ($1.32$--$1.42$), whereas CSRNet's rises from $1.20$ at
$49\times33$ to between $2.0$ and $2.2$ from $193\times130$ upward.

Run-to-run spread is larger for MobileCount at every resolution except
$49\times33$. Relative to the mean, CSRNet's spread narrows overall along the
ladder, from $12\%$ of its MAE at $49\times33$ to $0.8\%$ at $1544\times1038$,
though not monotonically: it rebounds to $4\%$ at $772\times519$. MobileCount's
shows no such trend, holding between $7\%$ and $18\%$ throughout. The best-epoch
ranges show a matching asymmetry: MobileCount's shift markedly earlier at the
top two resolutions ($100$--$252$ and $28$--$255$) relative to $357$--$579$ and
$367$--$693$ at $193\times130$ and $386\times260$, whereas CSRNet's remain late
at every resolution above the coarsest.

\begin{figure}
    \centering
    \includegraphics[width=0.96\linewidth]{images/mae_vs_memory_repeats.png}
    \caption{Accuracy--memory frontier: test MAE averaged over the $5$
    repeats against peak inference memory, for all twelve
    network~$\times$~resolution cells.}
    \label{fig:acc_vs_mem}
\end{figure}
Figure~\ref{fig:acc_vs_mem} plots test MAE against peak inference memory for all
twelve configurations. MobileCount is between $5.5$ and $8.1\times$ lighter than
CSRNet at every resolution. Model weights account for a static $13.1$\,MB of
this for MobileCount, $4.9\times$ less than CSRNet's $63.8$\,MB, and dominate the
footprint at the coarse end: at $49\times33$ peak memory is $13.4$\,MB for
MobileCount against $85.9$\,MB for CSRNet. Above $386\times260$ --- where the two
require $25.7$\,MB and $164.8$\,MB --- activations dominate instead and scale by
roughly $4\times$ per step, matching the growth in pixel count, reaching
$204.0$\,MB and $1648.8$\,MB respectively at $1544\times1038$.

\section{Discussion}
\label{sec:discussion}

\subsection{Configuring the Two Architectures}
\label{subsec:disc_hyper}

AdamW is the better optimiser for both networks. CSRNet's ImageNet-pretrained,
batch-norm-free VGG-16 frontend~\cite{li2018csrnet} is eroded by Adam's coupled
weight decay wherever gradients are small, leaving Adam uncompetitive at every
rate tested; MobileCount, batch-normalised and trained from scratch, is largely
invariant to the scale of its weights and tolerates coupled decay far better.
The same architectural difference drives the batch-size split: MobileCount's
normalisation statistics degenerate when a batch holds a single sample, whereas
CSRNet contains no batch normalisation and is best in exactly that setting,
consistent with its original publication~\cite{li2018csrnet}. Neither
configuration transfers --- batch size $1$ costs MobileCount a factor of seven
in validation MAE, and MobileCount's learning rate is one at which CSRNet does
not train at all. Hyperparameters published alongside a crowd-counting
architecture are therefore not portable to a different backbone even with task
and dataset fixed, and the framework defaults are a poor start for either.

\subsection{Accuracy Scaling with Input Fidelity (RQ1)}
\label{subsec:disc_accuracy}

\paragraph{What bounds accuracy changes along the ladder.}
CSRNet's collapse at the coarsest resolution is explained by its prediction grid
rather than its capacity: at $H/8$ it has fewer output cells than larvae to
represent, so no arrangement of weights can express the target and the reported
count comes from a near-uniform mass (Figure~\ref{fig:density_maps}).
MobileCount, predicting at $H/4$, is never grid-limited on this ladder, which
also explains why the two networks' largest and smallest gains fall on the same
step --- the one where CSRNet's grid first exceeds the mean count. Through the
middle of the ladder the inputs carry enough signal for both networks but too
little fine detail for extra capacity to be convertible; accuracy there is
bounded by input fidelity, and only at the top by architecture.

\paragraph{Two error distributions, not one ranking.}
The metrics disagree because the networks fail differently. CSRNet's rising
MSE-to-MAE ratio suggests its capacity buys a better typical frame at the cost
of a heavier tail of severe miscounts, while MobileCount's flat ratio indicates
errors tightly distributed about its mean. Since rearing decisions are threshold
judgements on single observations, the flatter distribution is the more useful
property, partly offsetting CSRNet's MAE advantage.

\paragraph{An optimisation caveat.}
MobileCount's early best checkpoints at the finest resolutions and its
degradation on both splits point to a learning rate transported from
$386\times260$ being too aggressive at sixteen times the pixel count. Its
high-resolution results are plausibly a lower bound, and the gap to CSRNet at
the top of the ladder correspondingly overstated.

\subsection{Memory Scaling with Input Fidelity (RQ2)}
\label{subsec:disc_memory}

Peak inference memory has two regimes, and the transition between them
determines where downscaling stops paying. 

At the three smallest resolutions the
footprint is weight-dominated and effectively flat, so reducing resolution further saves almost nothing; the step to $386\times260$ is where activations catch the weights, and above it they dominate and peak usage tracks pixel count. Beyond $386\times260$ activations dominate and peak usage tracks pixel count. Downscaling below $193\times130$ therefore sacrifices substantial accuracy for negligible memory --- a floor on the useful
part of the ladder imposed from the memory side rather than the accuracy side.
Because MobileCount carries roughly a fifth of CSRNet's weights, its flat regime
sits an order of magnitude lower and the ratio between the two stays within
$5.5$--$8.1\times$ throughout. Backbone and resolution thus act on memory almost
independently: there is no input size at which the heavier network becomes the
cheaper one to run, and no downscaling that closes the gap.

\subsection{The Accuracy--Memory Frontier (RQ3)}
\label{subsec:disc_pareto}

All six MobileCount configurations are Pareto-optimal, whereas four of CSRNet's
six are dominated outright. The decisive cell is MobileCount at $386\times260$:
it matches CSRNet at the same resolution to within $0.001$ MAE --- far below
either network's run-to-run spread --- for $25.7$\,MB against $164.8$\,MB, and
is both more accurate and cheaper than all three of CSRNet's coarser
configurations. CSRNet enters the frontier only at the high-accuracy end, buying
an $11\%$ MAE reduction over MobileCount's best for $2.3\times$ the memory at
$772\times519$, and $24\%$ for $8.1\times$ the memory at $1544\times1038$. Both purchases are on the MAE metric alone: on MSE
MobileCount remains ahead at both of those resolutions, so CSRNet's corner of
the frontier is narrower than Figure~\ref{fig:acc_vs_mem} alone suggests.

\subsection{Deployment on the Jetson Nano (RQ4)}
\label{subsec:disc_deployment}

The Jetson Nano shares a nominal $4$\,GB of RAM between CPU and GPU, of which
about $3.9$\,GB is addressable and a further ${\approx}750$\,MB is consumed by
the operating system, drivers, and CUDA runtime. Every one of the twelve
configurations fits the remaining ${\approx}3.15$\,GB: the
largest footprint measured, CSRNet's $1.6$\,GB at $1544\times1038$, occupies
about half of it, and the smallest, MobileCount's $13.4$\,MB at $49\times33$,
about $0.4\%$. Memory is therefore not the binding constraint on this hardware,
and the question becomes how much accuracy the application actually requires.
Counting systems in operational insect rearing report errors of
$2.3$--$6\%$~\cite{mamai2019, benedict2020}; from $193\times130$ upward both
networks stay below $2\%$ of the dataset's mean count of $47.9$ larvae, so every
configuration from that resolution up already clears the operational bar, and
memory spent climbing further buys accuracy the application does not use.

Three considerations converge on MobileCount at $386\times260$. It costs
$9.5$\,MB more than $193\times130$, the cheapest cell clearing the accuracy bar,
and halves the counting error for that outlay; it sits at the crossover between the two memory regimes, the last resolution at which weights still account for half its footprint, before each further doubling multiplies memory by roughly four; and
it is the resolution at which the hyperparameters were tuned, hence the last
below the instability discussed in Section~\ref{subsec:disc_accuracy}. Where the memory budget is tighter than this,
$193\times130$ remains the smallest configuration that still clears the bar;
where average error must be pushed lower still, CSRNet at $772\times519$ is the
next Pareto-optimal step, at roughly $18\times$ the memory.

\section{Limitations}
\label{sec:limitations}

\paragraph{Tuning protocol}
Hyperparameters were tuned one axis at a time in a fixed order, each
configuration trained once and the learning rate swept at decade granularity.
This OFAT procedure cannot observe interactions between stages, and the
configuration it converges on is conditional on the order in which the axes were
visited. Because augmentation came first, both networks were still running the
default Adam configuration when the crop was selected, and the scheme was never
re-validated under AdamW. The $800$-epoch budget was also tight in the middle of
the ladder, where MobileCount's best checkpoint fell as late as epoch $693$.

\paragraph{Hyperparameters transported from a single resolution}
The central limitation is that hyperparameters were optimised at $386\times260$
and applied unchanged across the ladder. This is what allows differences across
resolutions to be read as effects of input scale, but it confounds the top of
the ladder: MobileCount's high-resolution instability cannot be separated from a
transported learning rate simply too aggressive at four to sixteen times the
pixel count. Had it trained stably at $772\times519$ and above, it might have
approached the heavier network's accuracy at a fraction of the memory.
Separating an optimisation artefact from a genuine capacity ceiling requires
per-resolution retuning, which the frozen protocol excludes by construction.

\paragraph{Memory measurement and deployment estimate}
Peak memory is the PyTorch allocator's peak allocation for a single
\texttt{fp32} forward pass on a desktop GPU, and so excludes batching, the
preprocessing pipeline and camera capture buffers. Quantisation or a TensorRT
deployment would act mainly on the weight footprint that dominates the smallest
resolutions, and CSRNet, being weight-heavy, has proportionally more to gain, so
the $5.5$--$8.1\times$ ratio is an upper bound on the gap under an optimised
runtime. The Jetson Nano figures are estimates against the module's nominal
capacity rather than measurements on the device.

\paragraph{Dataset and task scope}
All results come from a single dataset. TenebrioVision is the only public
benchmark reflecting industrial rearing densities, but it is captured under
controlled conditions with a fixed viewpoint and counts topping out at $100$
larvae per image, whereas production trays are denser and more variable in
substrate, lighting and larval instar. The figures reported here should be read
as an upper bound on production imagery, and the comparison against the
$2.3$--$6\%$ range reported for operational counting
systems~\cite{mamai2019, benedict2020} is across measurement settings rather
than like for like. Ground-truth points blurred with a single fixed $\sigma$ per
resolution likewise suit roughly uniform larva size but would not transfer to
trays of mixed instars.

\section{Conclusion}
% ============================================================
This thesis compared two strategies for density-based counting of
\textit{Tenebrio molitor} larvae under edge-hardware memory constraints: a
heavyweight network on downscaled inputs against a lightweight one with headroom
for larger inputs. Density estimation proved viable for the task and remarkably
robust to downscaling, both networks counting to within one larva at
$193\times130$, a sixteenth of the native image in each dimension.

Beyond that point the two levers behave differently. Accuracy keeps improving
with resolution but with sharply diminishing returns above $386\times260$,
whereas memory is flat below that point and proportional to pixel count above
it, the two backbones separated by a near-constant $5.5$--$8.1\times$. The
answer to the main research question is therefore that input fidelity, not
architectural capacity, is where a memory budget should go---but only as far as
the knee, since MobileCount occupies the entire accuracy--memory frontier apart
from its high-accuracy corner.

On the Jetson Nano memory turns out not to be the binding constraint at all:
every configuration fits, and every one from $193\times130$ upward already meets
the error tolerances reported for operational insect rearing. MobileCount at
$386\times260$ is the recommended operating point, counting to within $1\%$ of
the mean tray population at $25.7$\,MB. The heavyweight network is justified
only where sub-$0.3$-larva mean error is a hard requirement, and even that niche
may be contestable, since MobileCount's high-resolution ceiling reflects
optimisation instability under transported hyperparameters rather than a
demonstrated capacity limit. Confirming this on the device, with per-resolution
retuning, is the natural next step.
\clearpage
\onecolumn
\section*{Appendix}
\label{sec:appendix}

\begin{figure}[h]
    \centering
    \includegraphics[width=\textwidth, height=0.8\textheight, keepaspectratio]{images/density_example_80_18_allres.png}
    \caption{Predicted density maps across the resolution ladder for test image
    \texttt{80\_18} (ground-truth count $80$). Columns show the input, the
    ground-truth density map, and the CSRNet and MobileCount predictions;
    all density maps share the colour scale at right.}
    \label{fig:density_maps}
\end{figure}

\begin{table*}[t]
    \centering
    \caption{Cross-resolution test results. Values are mean $\pm$ standard
    deviation over five runs; ``Epoch'' gives the range of best-validation
    epochs across those runs. Best value per column in bold.
    $^{*}$Resolution at which hyperparameters were tuned.}
    \label{tab:resolution_test}
    \small
    \begin{tabular}{l c c c c c c}
        \toprule
        & \multicolumn{3}{c}{CSRNet} & \multicolumn{3}{c}{MobileCount} \\
        \cmidrule(lr){2-4} \cmidrule(lr){5-7}
        Resolution & Epoch & MAE & MSE & Epoch & MAE & MSE \\
        \midrule
        $49\times33$      & $33$         & $6.551 \pm 0.782$          & $7.849 \pm 0.689$          & $196$--$361$ & $2.711 \pm 0.478$          & $3.584 \pm 0.670$ \\
        $97\times65$      & $145$--$328$ & $1.395 \pm 0.131$          & $2.307 \pm 0.227$          & $43$--$309$  & $2.549 \pm 0.258$          & $3.543 \pm 0.290$ \\
        $193\times130$    & $231$--$351$ & $0.899 \pm 0.046$          & $1.866 \pm 0.114$          & $357$--$579$ & $0.905 \pm 0.084$          & $1.236 \pm 0.131$ \\
        $386\times260^{*}$& $349$--$425$ & $0.472 \pm 0.005$          & $0.992 \pm 0.009$          & $367$--$693$ & $0.471 \pm 0.034$          & $0.658 \pm 0.041$ \\
        $772\times519$    & $145$--$536$ & $0.279 \pm 0.012$          & $0.614 \pm 0.047$          & $100$--$252$ & $0.429 \pm 0.045$          & $0.582 \pm 0.042$ \\
        $1544\times1038$  & $229$--$268$ & $\mathbf{0.238 \pm 0.002}$ & $\mathbf{0.483 \pm 0.012}$ & $28$--$255$  & $\mathbf{0.312 \pm 0.035}$ & $\mathbf{0.444 \pm 0.039}$ \\
        \bottomrule
    \end{tabular}
\end{table*}

\begin{table*}[t]
    \centering
    \caption{E4 --- Train- and validation-split MAE at the best-validation
    checkpoint and at the end of the fixed $1500$-epoch budget, for each
    resolution. Best validation MAE per architecture in bold.}
    \label{tab:e4_budget}
    \small
    \setlength{\tabcolsep}{4pt}
    \begin{tabular}{@{}l rr rr rr rr@{}}
        \toprule
        & \multicolumn{4}{c}{CSRNet} & \multicolumn{4}{c}{MobileCount} \\
        \cmidrule(lr){2-5} \cmidrule(lr){6-9}
        & \multicolumn{2}{c}{Checkpoint} & \multicolumn{2}{c}{Epoch 1500}
        & \multicolumn{2}{c}{Checkpoint} & \multicolumn{2}{c}{Epoch 1500} \\
        \cmidrule(lr){2-3} \cmidrule(lr){4-5} \cmidrule(lr){6-7} \cmidrule(lr){8-9}
        Resolution & Train & Val & Train & Val & Train & Val & Train & Val \\
        \midrule
        $49\times33$      & $5.94$ & $5.96$ & $12.63$ & $12.52$ & $2.61$ & $2.68$ & $6.29$ & $6.12$ \\
        $97\times65$      & $0.81$ & $1.03$ & $1.38$  & $1.46$  & $2.06$ & $1.91$ & $2.39$ & $2.49$ \\
        $193\times130$    & $0.29$ & $0.53$ & $0.21$  & $0.60$  & $0.86$ & $0.86$ & $1.15$ & $1.39$ \\
        $386\times260$    & $0.16$ & $0.29$ & $0.07$  & $0.32$  & $0.40$ & $\mathbf{0.39}$ & $0.46$ & $0.49$ \\
        $772\times519$    & $0.08$ & $0.19$ & $0.04$  & $0.23$  & $0.44$ & $0.43$ & $1.46$ & $1.59$ \\
        $1544\times1038$  & $0.09$ & $\mathbf{0.18}$ & $0.04$ & $0.24$ & $0.33$ & $0.43$ & $1.49$ & $1.58$ \\
        \bottomrule
    \end{tabular}
\end{table*}

\clearpage
\twocolumn

\bibliographystyle{IEEEtran}
\bibliography{references}
\end{document}